% ===== PRÉAMBULE CANONIQUE — à coller VERBATIM en tête de CHAQUE .tex =====
\documentclass[11pt,a4paper]{article}
\usepackage[utf8]{inputenc}\usepackage[T1]{fontenc}\usepackage{lmodern}
\usepackage[french]{babel}
\usepackage{amsmath,amssymb,mathtools}\usepackage{icomma}
\usepackage{siunitx}\sisetup{locale=FR}  % \qty{}{}, \si{}, \ang{}
\usepackage{xcolor}\usepackage{colortbl}
\usepackage{tikz}\usepackage{pgfplots}\pgfplotsset{compat=1.18}
\usepackage[siunitx,europeanresistors]{circuitikz} % circuits (élec.)
\usepackage[most]{tcolorbox}\usepackage{enumitem}\usepackage{array,booktabs}
\usepackage[a4paper,margin=2.2cm]{geometry}
\usetikzlibrary{arrows.meta,calc,angles,quotes,positioning,patterns,%
  backgrounds,shapes.geometric,decorations.pathmorphing,decorations.markings,intersections}
\definecolor{primary}{RGB}{222,49,99}\definecolor{primarydark}{RGB}{150,22,55}
\definecolor{defcol}{RGB}{0,114,178}\definecolor{methcol}{RGB}{52,92,148}
\definecolor{excol}{RGB}{0,135,90}\definecolor{attncol}{RGB}{200,40,55}
\tcbset{boxbase/.style={enhanced,breakable,boxrule=0.9pt,arc=2.5pt,
  left=3mm,right=3mm,top=2.3mm,bottom=2.3mm,fonttitle=\bfseries,coltitle=white}}
% --- boîtes TUTORIEL (noms EXACTS, ne pas renommer) ---
\newtcolorbox{cle}[1][]{boxbase,colback=defcol!6,colframe=defcol,
  title={\textsc{À retenir}\ifx\relax#1\relax\relax\else\textmd{ : \itshape #1}\fi}}
\newtcolorbox{methode}[1][]{boxbase,colback=methcol!7,colframe=methcol,sharp corners,
  title={\textsc{Méthode}\ifx\relax#1\relax\relax\else\textmd{ --- \itshape #1}\fi}}
\newtcolorbox{exemplebox}[1][]{boxbase,colback=excol!5,colframe=excol,
  title={\textsc{Exemple}\ifx\relax#1\relax\relax\else\textmd{ : \itshape #1}\fi}}
\newtcolorbox{piege}[1][]{boxbase,colback=attncol!6,colframe=attncol,
  title={\textsc{Piège}\ifx\relax#1\relax\relax\else\textmd{ : \itshape #1}\fi}}
\newtcolorbox{remarque}[1][]{boxbase,colback=black!4,colframe=black!45,
  title={\textsc{Remarque}\ifx\relax#1\relax\relax\else\textmd{ : \itshape #1}\fi}}
% --- boîtes EXERCICE / CORRIGÉ (noms EXACTS, auto-comptées) ---
\newtcolorbox[auto counter]{exo}[1][]{boxbase,colback=primary!4,colframe=primary,
  title={\textsc{Exercice}~\thetcbcounter\ifx\relax#1\relax\relax\else\textmd{ --- \itshape #1}\fi}}
\newtcolorbox[auto counter]{corrige}[1][]{boxbase,colback=excol!4,colframe=excol,
  title={\textsc{Corrigé}~\thetcbcounter\ifx\relax#1\relax\relax\else\textmd{ --- \itshape #1}\fi}}
\newcommand{\vv}[1]{\vec{#1}}\newcommand{\ds}{\displaystyle}
\newcommand{\R}{\mathbb{R}}\newcommand{\N}{\mathbb{N}}\newcommand{\Z}{\mathbb{Z}}
% (un \usepackage{fancyhdr}+en-tête est possible ; il est ignoré côté web)
% =========================================================================
\begin{document}
\sisetup{per-mode=symbol}

\begin{exo}[le choix de la référence est arbitraire]
Un drone de masse $m=\qty{4}{\kilogram}$ est immobile à $\qty{3}{\metre}$ au-dessus du sol. Une table
mesure $\qty{1}{\metre}$ de haut. On prend $g=\qty{10}{\metre\per\second\squared}$.
\begin{enumerate}[label=\alph*)]
  \item Calcule $E_p$ du drone en prenant la référence ($E_p=0$) au \emph{sol}.
  \item Calcule $E_p$ en prenant la référence sur le \emph{dessus de la table}.
  \item Le drone descend et se pose sur la table. Calcule $\Delta E_p$ dans \emph{chacune} des deux
        références. Que constates-tu ?
\end{enumerate}
\end{exo}

\begin{exo}[l'énergie élastique varie comme le carré de l'élongation]
Un ressort a une raideur $k=\qty{300}{\newton\per\metre}$.
\begin{enumerate}[label=\alph*)]
  \item Calcule $E_p$ pour une élongation $x_1=\qty{4}{\centi\metre}$.
  \item Calcule $E_p$ pour une élongation $x_2=\qty{8}{\centi\metre}$. Par quel facteur $E_p$ a-t-elle
        été multipliée quand l'élongation a doublé ?
\end{enumerate}
\end{exo}

\begin{exo}[énergie stockée entre deux élongations]
Un ressort de raideur $k=\qty{500}{\newton\per\metre}$ est déjà étiré de $x_1=\qty{2}{\centi\metre}$. On
l'étire davantage jusqu'à $x_2=\qty{6}{\centi\metre}$. Le travail supplémentaire à fournir vaut
$W=\tfrac12 k\,x_2^2-\tfrac12 k\,x_1^2$. Calcule ce travail (pense à convertir en mètres).
\end{exo}

\begin{exo}[influence de la masse et de la raideur sur la période]
La période d'un ressort est $T=2\pi\sqrt{m/k}$.
\begin{enumerate}[label=\alph*)]
  \item Par quel facteur la période est-elle multipliée si l'on \emph{quadruple} la masse $m$
        (à $k$ constant) ? Justifie avec la formule.
  \item Par quel facteur la période est-elle multipliée si l'on \emph{quadruple} la raideur $k$
        (à $m$ constant) ?
\end{enumerate}
\end{exo}

\begin{exo}[l'énergie potentielle ne dépend que de la hauteur]
On élève une caisse de masse $m=\qty{20}{\kilogram}$ jusqu'à une hauteur $h=\qty{3}{\metre}$, une fois
en la hissant verticalement, une autre fois en la poussant le long d'une rampe de longueur
$\qty{6}{\metre}$. On prend $g=\qty{10}{\metre\per\second\squared}$.
\begin{enumerate}[label=\alph*)]
  \item Calcule la variation d'énergie potentielle $\Delta E_p$ pour atteindre $h=\qty{3}{\metre}$.
  \item Cette variation dépend-elle du chemin suivi (rampe ou montée verticale) ? Justifie.
\end{enumerate}
\end{exo}
\end{document}
